\begin{longtable}[l]{ |m{5.2cm}|c|c|c|c|m{5.5cm}| }    
\caption{BLOCK0 参数}
    \label{table:BLOCK0_Parameter}
    \\\hline
    \rowcolor{lightgray}
    参数 & \thead{偏移} & \thead{位宽} & \thead{硬件使用} & \thead{\hyperref[fielddesc:EFUSEWRDIS]{EFUSE\_WR\_DIS} \\烧写保护位}
    & 描述 \\
    \endfirsthead
    \hline
    \rowcolor{lightgray}
    参数 & \thead{偏移} & \thead{位宽} & \thead{硬件使用} & \thead{\hyperref[fielddesc:EFUSEWRDIS]{EFUSE\_WR\_DIS} \\烧写保护位}
    & 描述 \\
    \endhead    
    \endfoot
    \endlastfoot
    \hline   
    \hyperref[fielddesc:EFUSEWRDIS]{\textbf{EFUSE\_WR\_DIS}} & 0 & 32 & Y & N/A & 表示是否禁止 eFuse 烧写 \\
    \hline
    \hyperref[fielddesc:EFUSERDDIS]{\textbf{EFUSE\_RD\_DIS}} & 32 & 7 & Y & 0  & 表示是否禁止用户读取 eFuse BLOCK4 \verb+~+ 10 的内容 \\
    \hline
    \hyperref[fielddesc:EFUSEDISICACHE]{EFUSE\_DIS\_ICACHE} & 40 &  1 & Y & 2  & 表示是否关闭 ICache \\
    \hline
    \hyperref[fielddesc:EFUSEDISDCACHE]{EFUSE\_DIS\_DCACHE} & 41 & 1 & Y & 2  & 表示是否关闭 DCache \\
    \hline
    \hyperref[fielddesc:EFUSEDISDOWNLOADICACHE]{EFUSE\_DIS\_DOWNLOAD\_ ICACHE} & 42 & 1 & Y & 2 & 表示是否在 Download 模式下关闭 ICache\\   
    \hline
    \hyperref[fielddesc:EFUSEDISDOWNLOADDCACHE]{EFUSE\_DIS\_DOWNLOAD\_ DCACHE} & 43 & 1 & Y & 2 & 表示是否在 Download 模式下关闭 DCache\\      
    \hline
    \hyperref[fielddesc:EFUSEDISFORCEDOWNLOAD]{EFUSE\_DIS\_FORCE\_DOWNLOAD} & 44 & 1 & Y & 2 & 表示是否禁止强制芯片进入 Download 模式\\   
    \hline
    \hyperref[fielddesc:EFUSEDISUSB]{EFUSE\_DIS\_USB} & 45 &1 & Y & 2 & 表示是否关闭 USB OTG 功能 \\     
    \hline
    \hyperref[fielddesc:EFUSEDISCAN]{EFUSE\_DIS\_CAN} & 46 & 1 & Y & 2 & 表示是否关闭 TWAI 控制器功能 \\  
    \hline
    \hyperref[fielddesc:EFUSEDISBOOTREMAP]{EFUSE\_DIS\_BOOT\_REMAP} & 47 & 1 & Y & 2 & 表示 RAM 空间是否可以映射到 ROM 空间 \\   
    \hline
    \hyperref[fielddesc:EFUSESOFTDISJTAG]{EFUSE\_SOFT\_DIS\_JTAG} & 49 & 1 & Y & 2 & 表示是否软禁用 JTAG \\    
    \hline
    \hyperref[fielddesc:EFUSEHARDDISJTAG]{EFUSE\_HARD\_DIS\_JTAG} & 50 & 1 & Y & 2 & 表示是否硬禁用 JTAG（永久） \\   
    \hline
    \hyperref[fielddesc:EFUSEDISDOWNLOADMANUALENCRYPT]{EFUSE\_DIS\_DOWNLOAD\_ MANUAL\_ENCRYPT} & 51 &  1 & Y & 2 & 表示是否在 download boot 模式下禁用 flash 加密功能 \\ 
    \hline
    \hyperref[fielddesc:EFUSEUSBEXCHGPINS]{EFUSE\_USB\_EXCHG\_PINS} & 56 & 1 & Y & 30 & 表示是否交换 USB D+/D- 管脚 \\
    \hline
    \hyperref[fielddesc:EFUSEEXTPHYENABLE]{EFUSE\_EXT\_PHY\_ENABLE} & 57 & 1 & N & 30 & 表示是否使能外部 USB PHY \\
    \hline
    \hyperref[fielddesc:EFUSEUSBFORCENOPERSIST]{EFUSE\_USB\_FORCE\_NOPERSIST} & 58 & 1 & N & 30 & 表示是否强制将 USB BVALID 设置为 1 \\    
    \hline
    \hyperref[fielddesc:EFUSEVDDSPIXPD]{EFUSE\_VDD\_SPI\_XPD} & 68 & 1 & Y & 3 & 表示 VDD\_SPI\_FORCE 为 1 时，VDD\_SPI 调节器是否上电 \\        
    \hline
    \hyperref[fielddesc:EFUSEVDDSPITIEH]{EFUSE\_VDD\_SPI\_TIEH} & 69 & 1 & Y & 3 & 表示 VDD\_SPI\_FORCE 为 1 时，VDD\_SPI 的电压选择 \\   
    \hline
    \hyperref[fielddesc:EFUSEVDDSPIFORCE]{EFUSE\_VDD\_SPI\_FORCE} & 70 & 1 & Y & 3 & 表示是否使用 XPD\_VDD\_PSI\_REG 和 VDD\_SPI\_TIEH 配置 VDD\_SPI LDO \\ 
    \hline
    \hyperref[fielddesc:EFUSEWDTDELAYSEL]{EFUSE\_WDT\_DELAY\_SEL} & 80 & 2 & Y & 3 & 表示是否选择 RTC 看门狗超时阈值 \\
    \hline
    \hyperref[fielddesc:EFUSESPIBOOTCRYPTCNT]{EFUSE\_SPI\_BOOT\_CRYPT\_CNT} & 82 & 3 & Y & 4 & 表示是否使能 SPI boot 加解密 \\
    \hline
    \hyperref[fielddesc:EFUSESECUREBOOTKEYREVOKE0]{EFUSE\_SECURE\_BOOT\_KEY\_ REVOKE0} & 85 & 1 & N & 5 & 表示是否撤销第一个安全启动 (Secure Boot) 密钥 \\
    \hline
    \hyperref[fielddesc:EFUSESECUREBOOTKEYREVOKE1]{EFUSE\_SECURE\_BOOT\_KEY\_ REVOKE1} & 86 & 1 & N & 6 & 表示是否撤销第二个安全启动密钥 \\
    \hline
    \hyperref[fielddesc:EFUSESECUREBOOTKEYREVOKE2]{EFUSE\_SECURE\_BOOT\_KEY\_ REVOKE2} & 87 & 1 & N & 7 & 表示是否撤销第三个安全启动密钥 \\
    \hline
    \hyperref[fielddesc:EFUSEKEYPURPOSE0]{EFUSE\_KEY\_PURPOSE\_0} & 88 & 4 & Y & 8 & 表示 Key0 用途 (purpose)，见表 \ref{table: key_purpose} \\
    \hline
    \hyperref[fielddesc:EFUSEKEYPURPOSE1]{EFUSE\_KEY\_PURPOSE\_1} & 92 & 4 & Y & 9 & 表示 Key1 用途，见表 \ref{table: key_purpose} \\
    \hline
    \hyperref[fielddesc:EFUSEKEYPURPOSE2]{EFUSE\_KEY\_PURPOSE\_2} & 96 & 4 & Y & 10 & 表示 Key2 用途，见表 \ref{table: key_purpose} \\ 
    \hline
    \hyperref[fielddesc:EFUSEKEYPURPOSE3]{EFUSE\_KEY\_PURPOSE\_3} & 100 & 4 & Y & 11 &  表示 Key3 用途，见表 \ref{table: key_purpose} \\    
    \hline
    \hyperref[fielddesc:EFUSEKEYPURPOSE4]{EFUSE\_KEY\_PURPOSE\_4} & 104 & 4 & Y & 12 & 表示 Key4 用途，见表 \ref{table: key_purpose} \\    
    \hline
    \hyperref[fielddesc:EFUSEKEYPURPOSE5]{EFUSE\_KEY\_PURPOSE\_5} & 108 & 4 & Y & 13 & 表示 Key5 用途，见表 \ref{table: key_purpose} \\  
    \hline
    \hyperref[fielddesc:EFUSESECUREBOOTEN]{EFUSE\_SECURE\_BOOT\_EN} & 116 & 1 & N & 15 & 表示是否使能安全启动 \\ 
    \hline
    \hyperref[fielddesc:EFUSESECUREBOOTAGGRESSIVEREVOKE]{EFUSE\_SECURE\_BOOT\_ AGGRESSIVE\_REVOKE} & 117 & 1 & N & 16 & 表示是否使安全启动的撤销采用激进策略 \\
    \hline
    \hyperref[fielddesc:EFUSEFLASHTPUW]{EFUSE\_FLASH\_TPUW} & 124 & 4 & N & 18 & 表示上电后 flash 等待时间 \\    
    \hline    
    \hyperref[fielddesc:EFUSEDISDOWNLOADMODE]{EFUSE\_DIS\_DOWNLOAD\_MODE} & 128 & 1 & N & 18 & 表示是否关闭所有 download boot 模式\\       
    \hline
    \hyperref[fielddesc:EFUSEDISLEGACYSPIBOOT]{EFUSE\_DIS\_LEGACY\_SPI\_BOOT} & 129 & 1 & N & 18 & 表示是否关闭 Legacy SPI boot 模式 \\
    \hline
    \hyperref[fielddesc:EFUSEUARTPRINTCHANNEL]{EFUSE\_UART\_PRINT\_CHANNEL} & 130 & 1 & N & 18 & 表示打印 boot 信息的 UART 通道选择 \\
    \hline 
    \hyperref[fielddesc:EFUSEDISUSBDOWNLOADMODE]{EFUSE\_DIS\_USB\_DOWNLOAD\_ MODE} & 132 & 1 & N & 18 & 表示是否在 UART download boot 模式下关闭 USB OTG 下载功能 \\
    \hline 
    \hyperref[fielddesc:EFUSEENABLESECURITYDOWNLOAD]{EFUSE\_ENABLE\_SECURITY\_ DOWNLOAD} & 133 & 1 & N & 18 & 表示是否使能 UART 安全下载模式 \\
    \hline
    \hyperref[fielddesc:EFUSEUARTPRINTCONTROL]{EFUSE\_UART\_PRINT\_CONTROL} & 134 & 2 & N & 18 & 表示 UART boot 信息的打印方式 \\
    \hline
    \hyperref[fielddesc:EFUSEPINPOWERSELECTION]{EFUSE\_PIN\_POWER\_SELECTION} & 136 & 1 & N & 18 & 表示 GPIO33 \verb+~+ GPIO37 的电源选择 \\
    \hline
    \hyperref[fielddesc:EFUSEFLASHTYPE]{EFUSE\_FLASH\_TYPE} & 137 & 1 & N & 18 & 表示 flash 类型 \\
    \hline
    \hyperref[fielddesc:EFUSEFORCESENDRESUME]{EFUSE\_FORCE\_SEND\_RESUME} & 138 & 1 & N & 18 & 表示是否强制 ROM 代码在 SPI 启动过程中发送恢复指令 \\  
    \hline
    \hyperref[fielddesc:EFUSESECUREVERSION]{EFUSE\_SECURE\_VERSION} & 139 & 16 & N & 18 & 表示 IDF 安全版本 \\\hline
\end{longtable}
